\section{Supplementary Theory and Method Details}
\label{sec:si-theory-details}

% ======================================================================
\subsection{Differentiable Kohn--Sham SCF}
\label{sec:scf}

The Kohn--Sham (KS) equations for a restricted closed-shell system
are
\begin{equation}
  \hat{f}_\theta[\rho] \phi_i(\mathbf{r})
  =
  \left[-\frac{1}{2}\nabla^2 + v_{\rm ext}(\mathbf{r})
  + v_J[\rho](\mathbf{r}) + v_{\rm xc}^{\theta}[\rho](\mathbf{r})
  \right]\phi_i(\mathbf{r})
  = \epsilon_i \phi_i(\mathbf{r}),
  \label{eq:ks}
\end{equation}
where $\theta$ denotes the neural XC parameters,
$v_{\rm xc}^{\theta}(\mathbf{r}) = \delta E_{\rm xc}^{\theta}/
\delta\rho(\mathbf{r})$, and
$\rho(\mathbf{r}) = 2\sum_i n_i |\phi_i(\mathbf{r})|^2$ with
orbital occupations $n_i$ normalized to one for a doubly occupied
restricted orbital. In the atomic orbital (AO) basis
$\{\chi_\mu\}$,
\begin{equation}
  \mathbf{F}_\theta[\mathbf{D}]\,\mathbf{C}
  =
  \mathbf{S}\mathbf{C}\boldsymbol{\epsilon},
  \label{eq:generalized-eig}
\end{equation}
with
\begin{equation}
  \mathbf{F}_\theta[\mathbf{D}]
  = \mathbf{h} + \mathbf{J}[\mathbf{D}]
  - \alpha_\theta \mathbf{K}[\mathbf{D}]
  + \mathbf{V}_{\rm xc}^{\theta}[\mathbf{D}]
  + \mathbf{V}_{\rm nl}^{\theta}[\mathbf{D}].
  \label{eq:fock}
\end{equation}
Here $\mathbf{h}$ is the one-electron Hamiltonian, $\mathbf{J}$ and
$\mathbf{K}$ are Coulomb and exact-exchange matrices, $\alpha_\theta$
is the scalar exchange fraction used by the nonlocal hybrid term, and
$\mathbf{V}_{\rm nl}^{\theta}$ denotes additional nonlocal
contributions generated by local-HF or PT2 neural channels when those
channels are active in the SCF energy.

GradTDDFT treats molecular construction as a fixed data layer. For a
given geometry and basis set, PySCF supplies
\begin{equation}
  \mathcal{I} =
  \{\mathbf{h},\mathbf{S},(\mu\nu|\lambda\sigma),
  \chi_\mu(\mathbf{r}_g),\nabla\chi_\mu(\mathbf{r}_g),
  w_g,\mathbf{r}_g\}_{g=1}^{N_g},
  \label{eq:fixed-data}
\end{equation}
and the differentiable JAX program updates only the density,
J/K contractions, XC potential, Fock matrix, orbital coefficients,
TDDFT response matrices, and losses. This distinction is important:
an ``integral backend'' may generate $\mathcal{I}$ on CPU or GPU,
but training-time gradients do not require differentiating through
the external integral engine.

The SCF iteration is written as the fixed-point map
\begin{equation}
  \mathbf{D}^{(n+1)}
  =
  \Phi_\theta(\mathbf{D}^{(n)};\mathcal{I}),
  \qquad
  \mathbf{R}_\theta(\mathbf{D})
  =
  \mathbf{D}-\Phi_\theta(\mathbf{D};\mathcal{I}),
  \label{eq:scf-fixed-point}
\end{equation}
where $\Phi_\theta$ denotes Fock construction, generalized
diagonalization, occupation assignment, and density reconstruction.
The converged density $\mathbf{D}^{*}_\theta$ satisfies
$\mathbf{R}_\theta(\mathbf{D}^{*}_\theta)=0$. GradTDDFT implements
restricted and unrestricted SCF solvers, with damping, level shifting,
and DIIS-style acceleration in the reference paths.\cite{DIIS}

\emph{Differentiable training modes.} Let
$\mathcal{O}_{\theta}(\mathbf{D};\mathcal{I})$ denote the
observable map built from the XC energy, XC potential, TDDFT
response matrices, excitation energies, and oscillator strengths
evaluated at density
$\mathbf{D}$. For a reference target $\mathbf{y}$, define a generic
loss
\begin{equation}
  \mathcal{L}(\theta,\mathbf{D})
  =
  \ell\!\left(
    \mathcal{O}_{\theta}(\mathbf{D};\mathcal{I}),\mathbf{y}
  \right).
  \label{eq:generic-loss}
\end{equation}
GradTDDFT uses three mathematically distinct training modes.

\textbf{Fixed-density training.} A reference density
$\mathbf{D}^{\rm ref}$ is treated as an external variational point:
\begin{equation}
  \mathcal{J}_{\rm fd}(\theta)
  =
  \mathcal{L}(\theta,\mathbf{D}^{\rm ref}),
  \qquad
  \frac{d\mathbf{D}^{\rm ref}}{d\theta}=0,
  \label{eq:fd-objective}
\end{equation}
and therefore
\begin{equation}
  \frac{d\mathcal{J}_{\rm fd}}{d\theta}
  =
  \left.
  \frac{\partial \mathcal{L}(\theta,\mathbf{D})}
       {\partial\theta}
  \right|_{\mathbf{D}=\mathbf{D}^{\rm ref}} .
  \label{eq:fd-gradient}
\end{equation}
This mode updates the neural XC energy, potential, and response
kernel at a fixed density while excluding the self-consistent density
response.

\textbf{Explicit self-consistent training.} Starting from an initial
density $\mathbf{D}^{(0)}$, a finite SCF trajectory is generated by
\begin{equation}
  \mathbf{D}^{(n+1)}
  =
  \Phi_{\theta}(\mathbf{D}^{(n)};\mathcal{I}),
  \qquad n=0,\ldots,N-1,
  \label{eq:unrolled-scf}
\end{equation}
and the training objective is
\begin{equation}
  \mathcal{J}_{\rm exp}^{(N)}(\theta)
  =
  \mathcal{L}\!\left(\theta,\mathbf{D}^{(N)}_{\theta}\right).
  \label{eq:explicit-objective}
\end{equation}
The gradient follows the differentiated finite trajectory. With
$\boldsymbol{\lambda}^{(N)}
=\partial\mathcal{L}/\partial\mathbf{D}^{(N)}$ and
\begin{equation}
  \boldsymbol{\lambda}^{(n)}
  =
  \left(
    \frac{\partial\Phi_{\theta}}
         {\partial\mathbf{D}^{(n)}}
  \right)^{T}
  \boldsymbol{\lambda}^{(n+1)},
  \qquad n=N-1,\ldots,0,
  \label{eq:explicit-adjoint}
\end{equation}
the parameter derivative is
\begin{equation}
  \frac{d\mathcal{J}_{\rm exp}^{(N)}}{d\theta}
  =
  \frac{\partial\mathcal{L}}{\partial\theta}
  +
  \sum_{n=0}^{N-1}
  \left(\boldsymbol{\lambda}^{(n+1)}\right)^{T}
  \frac{\partial\Phi_{\theta}(\mathbf{D}^{(n)};\mathcal{I})}
       {\partial\theta}.
  \label{eq:explicit-gradient}
\end{equation}
This mode preserves the dependence of the loss on every retained SCF
iterate.

\textbf{Implicit self-consistent training.} The SCF density is first
defined as the fixed point
\begin{equation}
  \mathbf{R}_{\theta}(\mathbf{D}^{*}_{\theta})=0,
  \qquad
  \mathbf{R}_{\theta}(\mathbf{D})
  =
  \mathbf{D}-\Phi_{\theta}(\mathbf{D};\mathcal{I}),
  \label{eq:implicit-fixed-point}
\end{equation}
and the objective is
\begin{equation}
  \mathcal{J}_{\rm imp}(\theta)
  =
  \mathcal{L}\!\left(\theta,\mathbf{D}^{*}_{\theta}\right).
  \label{eq:implicit-objective}
\end{equation}
The density response is represented by the adjoint variable
$\boldsymbol{\lambda}$ satisfying
\begin{equation}
  \left(
    \frac{\partial \mathbf{R}_{\theta}}
         {\partial \mathbf{D}}
  \right)^{T}
  \boldsymbol{\lambda}
  =
  -
  \frac{\partial\mathcal{L}}
       {\partial\mathbf{D}^{*}},
  \label{eq:implicit-adjoint}
\end{equation}
which gives
\begin{equation}
  \frac{d\mathcal{J}_{\rm imp}}{d\theta}
  =
  \frac{\partial\mathcal{L}}{\partial\theta}
  +
  \boldsymbol{\lambda}^{T}
  \frac{\partial\mathbf{R}_{\theta}}
       {\partial\theta}.
  \label{eq:implicit-scf-gradient}
\end{equation}
This mode differentiates the converged SCF solution without treating
the finite iteration history as part of the objective.

% ======================================================================
\subsection{Linear-Response TDDFT}
\label{sec:tddft}

Within the linear-response formalism, excitation energies
$\Omega_I$ and transition amplitudes are obtained by solving the
Casida eigenvalue equation:\cite{Casida1995}
\begin{equation}
  \begin{pmatrix}
  \mathbf{A} & \mathbf{B} \\
  \mathbf{B}^* & \mathbf{A}^*
  \end{pmatrix}
  \begin{pmatrix} \mathbf{X}_I \\ \mathbf{Y}_I \end{pmatrix}
  = \Omega_I
  \begin{pmatrix}
  \mathbf{1} & \mathbf{0} \\
  \mathbf{0} & -\mathbf{1}
  \end{pmatrix}
  \begin{pmatrix} \mathbf{X}_I \\ \mathbf{Y}_I \end{pmatrix},
  \label{eq:casida}
\end{equation}
where $\mathbf{X}_I$ and $\mathbf{Y}_I$ are the excitation and
de-excitation amplitudes in the occupied-virtual MO rotation basis.

For restricted closed-shell singlet excitations, the orbital Hessian
matrix elements are
\begin{equation}
  \begin{aligned}
  A_{ia,jb} &=
    (\epsilon_a - \epsilon_i)\delta_{ij}\delta_{ab}
    + 2(ia|jb) - \alpha_\theta(ij|ab)
    + K_{ia,jb}^{\rm xc,A}
    + K_{ia,jb}^{\rm nl,A}, \\
  B_{ia,jb} &=
    2(ia|bj) - \alpha_\theta(ib|aj)
    + K_{ia,jb}^{\rm xc,B}
    + K_{ia,jb}^{\rm nl,B},
  \end{aligned}
  \label{eq:ab-matrix}
\end{equation}
where $(pq|rs)$ are two-electron integrals in the MO basis,
$\alpha_\theta$ is the effective Hartree--Fock exchange fraction,
$K^{\rm xc}$ is the adiabatic semilocal/neural local response
contribution (Section~\ref{sec:fxc}), and $K^{\rm nl}$ collects
additional nonlocal response terms from neural hybrid channels when
they are enabled. Indices $i,j$ denote occupied MOs and $a,b$
virtual MOs. In the default semilocal response path,
$K^{\rm xc,A}=K^{\rm xc,B}$; nonlocal channels may contribute
different $A$ and $B$ matrices.

The Tamm--Dancoff form (TDA) sets $\mathbf{B}=0$,
reducing Eq.~\ref{eq:casida} to the Hermitian eigenvalue problem
$\mathbf{A}\mathbf{X}_I = \Omega_I \mathbf{X}_I$.

From the solution vectors, physical observables are computed:
\begin{align}
  \boldsymbol{\mu}_I &=
    \sum_{ia} (X_{ia,I} + Y_{ia,I})
    \boldsymbol{\mu}_{ia}, \label{eq:transdip} \\
  f_I &= \frac{2}{3}\Omega_I |\boldsymbol{\mu}_I|^2,
    \label{eq:oscstr}
\end{align}
where $\boldsymbol{\mu}_{ia} = \langle i|\mathbf{r}|a\rangle$ are
MO transition dipole integrals, and $f_I$ are oscillator strengths.

% ======================================================================
\subsection{Implicit Differentiation of TDDFT Roots}
\label{sec:tddft-implicit-roots}

For self-consistent excited-state training, the response matrices
depend on the neural parameters in two ways. First, the semilocal
response tensor and nonlocal channel coefficients depend explicitly
on $\theta$. Second, the response problem is evaluated at the
converged SCF density $\mathbf{D}^*_\theta$, whose derivative is
handled by the implicit SCF adjoint in
Eqs.~\ref{eq:implicit-adjoint}--\ref{eq:implicit-scf-gradient}.
The TDDFT root derivative is then taken with respect to the response
operator built at this converged state. The eigensolver derivative
follows implicit differentiation of a converged dominant eigensolver,
so the backward pass is defined by the fixed-point/eigenpair
conditions rather than by the individual Davidson iterations.
\cite{Xie2020DominantEigensolverAD}
This is the root-energy differential used in the implementation. The
Davidson iterations are used to obtain the forward Ritz amplitudes,
but the operator applications inside the iterative solve and the final
amplitudes are detached from the reverse pass. The differentiable
root energy is then reconstructed by applying the current response
operator to those fixed amplitudes.

For TDA, write the response operator as
\begin{equation}
  \mathbf{A}_{\theta}\mathbf{X}_I
  =
  \Omega_I \mathbf{X}_I .
  \label{eq:tda-root-si}
\end{equation}
The iterative eigensolver returns a converged Ritz vector
$\widehat{\mathbf{X}}_I$. In the reverse pass, GradTDDFT treats this
vector as fixed and reconstructs the differentiable excitation energy
from the Rayleigh quotient
\begin{equation}
  \Omega_I^{\rm TDA}(\theta)
  =
  \frac{
    \widehat{\mathbf{X}}_I^{T}
    \mathbf{A}_{\theta}
    \widehat{\mathbf{X}}_I
  }{
    \widehat{\mathbf{X}}_I^{T}
    \widehat{\mathbf{X}}_I
  }.
  \label{eq:tda-rayleigh-si}
\end{equation}
For a scalar loss depending on TDA roots, this gives the
root-energy contribution
\begin{equation}
  \frac{d\mathcal{L}}{d\theta}
  \supset
  \sum_I
  \frac{\partial\mathcal{L}}{\partial\Omega_I}
  \frac{
    \widehat{\mathbf{X}}_I^{T}
    \left(d\mathbf{A}_{\theta}/d\theta\right)
    \widehat{\mathbf{X}}_I
  }{
    \widehat{\mathbf{X}}_I^{T}
    \widehat{\mathbf{X}}_I
  },
  \label{eq:tda-root-gradient-si}
\end{equation}
where $d\mathbf{A}_{\theta}/d\theta$ includes both explicit neural
kernel derivatives and the derivative of the converged SCF state.

For full TDDFT, define
\begin{equation}
  \mathbf{H}_{\theta}
  =
  \begin{pmatrix}
  \mathbf{A}_{\theta} & \mathbf{B}_{\theta} \\
  \mathbf{B}_{\theta} & \mathbf{A}_{\theta}
  \end{pmatrix},
  \qquad
  \mathbf{S}
  =
  \begin{pmatrix}
  \mathbf{1} & \mathbf{0} \\
  \mathbf{0} & -\mathbf{1}
  \end{pmatrix},
  \qquad
  \mathbf{Z}_I=
  \begin{pmatrix}\mathbf{X}_I\\\mathbf{Y}_I\end{pmatrix}.
  \label{eq:casida-hs-si}
\end{equation}
The Casida equation is
$\mathbf{H}_{\theta}\mathbf{Z}_I
=\Omega_I\mathbf{S}\mathbf{Z}_I$. With converged Ritz amplitudes
$\widehat{\mathbf{X}}_I,\widehat{\mathbf{Y}}_I$ held fixed in the
reverse pass, the differentiable root is
\begin{equation}
\begin{aligned}
  \Omega_I^{\rm TDDFT}(\theta)
  &=
  \frac{
  \widehat{\mathbf{X}}_I^{T}
  \left(\mathbf{A}_{\theta}\widehat{\mathbf{X}}_I
       +\mathbf{B}_{\theta}\widehat{\mathbf{Y}}_I\right)
  +
  \widehat{\mathbf{Y}}_I^{T}
  \left(\mathbf{B}_{\theta}\widehat{\mathbf{X}}_I
       +\mathbf{A}_{\theta}\widehat{\mathbf{Y}}_I\right)
  }{
  \widehat{\mathbf{X}}_I^{T}\widehat{\mathbf{X}}_I
  -
  \widehat{\mathbf{Y}}_I^{T}\widehat{\mathbf{Y}}_I
  } .
\end{aligned}
  \label{eq:casida-rayleigh-si}
\end{equation}
Therefore
\begin{equation}
\begin{aligned}
  \frac{d\Omega_I^{\rm TDDFT}}{d\theta}
  &=
  \frac{1}{N_I}
  \Big[
  \widehat{\mathbf{X}}_I^{T}
  \left(d\mathbf{A}_{\theta}\widehat{\mathbf{X}}_I
       +d\mathbf{B}_{\theta}\widehat{\mathbf{Y}}_I\right) \\
  &\quad+
  \widehat{\mathbf{Y}}_I^{T}
  \left(d\mathbf{B}_{\theta}\widehat{\mathbf{X}}_I
       +d\mathbf{A}_{\theta}\widehat{\mathbf{Y}}_I\right)
  \Big],
  \qquad
  N_I=
  \widehat{\mathbf{X}}_I^{T}\widehat{\mathbf{X}}_I
  -
  \widehat{\mathbf{Y}}_I^{T}\widehat{\mathbf{Y}}_I .
\end{aligned}
  \label{eq:casida-root-gradient-si}
\end{equation}
Here $d\mathbf{A}_{\theta}$ and $d\mathbf{B}_{\theta}$ denote total
parameter differentials through the response matrix actions, including
the implicit SCF dependence. This is the mathematical form of the
implementation strategy: the Davidson iterations provide the forward
Ritz vectors, while the reverse-mode derivative follows the current
response operator rather than the eigensolver iteration history.
Amplitude-dependent observables, such as oscillator strengths, are
evaluated from the same converged roots and transition dipole
contractions used in the forward calculation. In the current training
path, these observables use the fixed forward amplitudes returned by
the eigensolver; the formulas above should therefore be interpreted
as root-energy derivatives rather than full eigenvector-response
derivatives for oscillator strengths.

% ======================================================================
\subsection{Energy-Consistent \texorpdfstring{$f_{\rm xc}$}{fxc} Kernel via JAX Automatic Differentiation}
\label{sec:fxc}

The exchange-correlation kernel is the second functional derivative
of the XC energy with respect to the density:
\begin{equation}
  f_{\rm xc}(\mathbf{r}, \mathbf{r}') =
    \frac{\delta^2 E_{\rm xc}}
         {\delta\rho(\mathbf{r})\delta\rho(\mathbf{r}')}.
  \label{eq:fxcfd}
\end{equation}

In most quantum chemistry packages, $f_{\rm xc}$ is implemented
through analytic derivatives of standard XC functional forms
or via finite differences, while automatic differentiation provides
a systematic route to higher-order density-functional response
derivatives.\cite{Ekstrom2010ADResponse} For neural XC functionals, where the
functional form is a neural network with thousands of parameters,
analytic derivatives are infeasible and finite differences introduce
numerical noise that compromises gradient-based optimization.

GradTDDFT addresses this through an \emph{energy-consistent local response
tensor}. The XC energy is expressed as an integral over a local
energy contribution:
\begin{equation}
  E_{\rm xc}^{\theta}
  =
  \sum_g w_g\,
  e_{\rm xc}^{\theta}\bigl(\mathbf{u}_g,\mathbf{q}_g\bigr),
  \label{eq:eloc}
\end{equation}
where $\mathbf{u}_g$ contains the differentiable density variables
and $\mathbf{q}_g$ contains auxiliary frozen or separately treated
features. For the current restricted MGGA response path,
\begin{equation}
  \mathbf{u}_g =
  \left[
    \rho_g,\,
    \partial_x\rho_g,\,
    \partial_y\rho_g,\,
    \partial_z\rho_g,\,
    \tau_g
  \right],
  \label{eq:response-vars}
\end{equation}
with the GGA path obtained by omitting $\tau_g$. The auxiliary
features $\mathbf{q}_g$ include local exact-exchange descriptors
and optional projected PT2 descriptors used by the coefficient
network.
For a conventional LDA, GGA, or meta-GGA functional, the same
response core is obtained by taking
$e_{\rm xc}^{\theta}=e_{\rm xc}^{F}$ for the selected functional
$F$. Thus differentiable evaluation is used not only for neural
functionals but also to compute the response kernel of traditional
functionals from the local energy-density expression. The
corresponding data flow is summarized in
Fig.~1 of the main text.

The local $f_{\rm xc}$ contribution is constructed by contracting
the grid Hessian with transition response features:
\begin{equation}
  K_{ia,jb}^{\rm xc}
  =
  2\sum_g w_g
  \sum_{xy}
  H_{xy,g}^{\theta}
  P_{x,g}^{ia}
  P_{y,g}^{jb},
  \label{eq:fxccontraction}
\end{equation}
where
\begin{align}
  H_{xy,g}^{\theta} &=
    \frac{\partial^2 e_{\rm xc}^{\theta}}
         {\partial u_{x,g}\, \partial u_{y,g}},
    \label{eq:localhessian} \\
  P_{x,g}^{ia} &=
    \frac{\partial u_{x,g}}{\partial \kappa_{ia}},
    \label{eq:transresponse}
\end{align}
and $\kappa_{ia}$ are occupied-virtual orbital rotation parameters.
For example, the density feature has
$P_{\rho,g}^{ia}=2\phi_i(\mathbf{r}_g)\phi_a(\mathbf{r}_g)$ in
the restricted singlet convention; gradient and kinetic-energy
features are obtained by differentiating the corresponding AO-grid
expressions with respect to the same orbital rotation.

Operationally, GradTDDFT treats the pointwise local energy as a
scalar map
\begin{equation}
  \mathcal{E}_\theta:\mathbf{u}_g\mapsto
  e_{\rm xc}^{\theta}(\mathbf{u}_g,\mathbf{q}_g),
  \label{eq:pointwise-map}
\end{equation}
and evaluates $H_{xy,g}^{\theta}$ as the Hessian
$\nabla_{\mathbf{u}}^2\mathcal{E}_\theta$ by automatic
differentiation. The first derivative of the same scalar map gives
the SCF potential components
\begin{equation}
  v_{x,g}^{\theta}
  =
  \frac{\partial \mathcal{E}_\theta}
       {\partial u_{x,g}},
  \label{eq:potential-components}
\end{equation}
so the energy, potential, and response tensor stay internally
consistent. This eliminates finite-difference errors and correctly
includes derivatives of the learned mixing coefficients with
respect to the density variables.

For hybrid functionals with HF exchange, the ordinary exact-exchange
kernel is treated as a nonlocal contribution:
\begin{equation}
  K_{ia,jb}^{\rm HF,A} = -\alpha_\theta(ij|ab), \qquad
  K_{ia,jb}^{\rm HF,B} = -\alpha_\theta(ib|aj).
  \label{eq:hfkernel}
\end{equation}
For the neural HF channel, $\alpha_\theta$ is obtained from the
density-weighted average of the learned local exchange field,
\begin{equation}
  \alpha_x^\theta
  =
  \frac{\sum_g w_g \rho_g c_x^\theta(\mathbf{z}_g)}
       {\sum_g w_g \rho_g}.
  \label{eq:hf-scalar-response-si}
\end{equation}
The local HF descriptor $h_g$ remains part of the coefficient input,
but it is held fixed inside the local Hessian. The orbital Hessian
therefore uses Eq.~\ref{eq:hfkernel} with
$\alpha_\theta=\alpha_x^\theta$, rather than a separate response
feature for the projected local HF field.

For the PT2 channel, GradTDDFT follows the CIS(D)-type
double-hybrid excited-state form
\cite{HeadGordonCISD1994,GrimmeNeeseDoubleHybrid2007}. The local
pair gauge $p_g$ and the learned coefficient field define a
density-weighted correlation coefficient
\begin{equation}
  a_c^\theta
  =
  \operatorname{clip}_{[0,1]}
  \left[
  \frac{\sum_g w_g \rho_g c_2^\theta(\mathbf{z}_g)}
       {\sum_g w_g \rho_g}
  \right].
  \label{eq:pt2-scalar-correlation-si}
\end{equation}
The linear-response problem is solved first with the semilocal
kernel and the scalar HF exchange kernel. The perturbative doubles
contribution is then added after the TDA or Casida eigenproblem, so
the PT2 local channel does not enter the response matrix as a local
second derivative. For each root,
\begin{equation}
  \Omega_I
  =
  \Omega_I^{\rm LR}
  + a_c^\theta \Delta_I^{\rm CIS(D)}.
  \label{eq:pt2-cisd-corrected-root-si}
\end{equation}
All quantities in $\Delta_I^{\rm CIS(D)}$ are built from the current
MO coefficients, orbital energies, two-electron integrals, and
singles amplitudes. The correction is root-specific and is evaluated
without spin-component scaling or damping in the present form.

Let $i,j,k$ denote occupied spin orbitals, $a,b,c$ virtual spin
orbitals, and
\begin{equation}
  \langle pq||rs\rangle=(pr|qs)-(ps|qr)
  \label{eq:cisd-antisym-integral-si}
\end{equation}
in chemists' notation. Define the positive orbital-energy gap
\begin{equation}
  D_{ij}^{ab}
  =
  \epsilon_a+\epsilon_b-\epsilon_i-\epsilon_j,
  \qquad
  t_{ij}^{ab}
  =
  -\frac{\langle ij||ab\rangle}{D_{ij}^{ab}},
  \label{eq:cisd-mp2-amplitudes-si}
\end{equation}
where $t_{ij}^{ab}$ is the corresponding MP2 doubles amplitude in
this denominator convention. Let $B_{ia}^{I}$ denote the singles
transition amplitude used by the correction, with $B=X$ for TDA and
$B=X+Y$ for full TDDFT. The direct CIS(D) contribution is
\begin{equation}
  \Delta_{I,{\rm dir}}^{\rm CIS(D)}
  =
  -\frac{1}{4}
  \sum_{ijab}
  \frac{\left(U_{ijab}^{I}\right)^2}
       {\epsilon_a+\epsilon_b-\epsilon_i-\epsilon_j-\Omega_I^{\rm LR}},
  \label{eq:cisd-direct-si}
\end{equation}
where
\begin{equation}
\begin{aligned}
  U_{ijab}^{I}
  &=
  \sum_c
  \left(\langle ab||cj\rangle B_{ic}^{I}
       -\langle ab||ci\rangle B_{jc}^{I}\right)  \\
  &\quad+
  \sum_k
  \left(\langle ka||ij\rangle B_{kb}^{I}
       -\langle kb||ij\rangle B_{ka}^{I}\right).
\end{aligned}
  \label{eq:cisd-u-vector-si}
\end{equation}
The denominator in Eq.~\ref{eq:cisd-direct-si} is
$D_{ij}^{ab}-\Omega_I^{\rm LR}$; the excitation energy shifts the
doubles manifold relative to the singles state.

The indirect term accounts for the coupling of the CIS singles vector
with ground-state MP2 amplitudes. The intermediates used in
GradTDDFT are
\begin{align}
  R_{ab}
  &=
  -\sum_{jkc}
  (jc|kb)\,t_{jk}^{ca},
  \label{eq:cisd-rab-si}
  \\
  R_{ij}
  &=
  -\sum_{kab}
  (ja|kb)\,t_{ik}^{ab},
  \label{eq:cisd-rij-si}
  \\
  W_{ia}^{I}
  &=
  \sum_{jkbc}
  \langle jk||bc\rangle\,
  t_{ik}^{ac}\,
  B_{jb}^{I}.
  \label{eq:cisd-wia-si}
\end{align}
The unscaled indirect correction is then
\begin{equation}
\begin{aligned}
  \Delta_{I,{\rm ind}}^{\rm CIS(D)}
  &=
  \sum_{iab}
  B_{ia}^{I}B_{ib}^{I}R_{ab}
  +
  \sum_{ijc}
  B_{ic}^{I}B_{jc}^{I}R_{ij}  \\
  &\quad+
  \sum_{ia}
  B_{ia}^{I}W_{ia}^{I}.
\end{aligned}
  \label{eq:cisd-indirect-si}
\end{equation}
The total second-order root correction is
\begin{equation}
  \Delta_I^{\rm CIS(D)}
  =
  \Delta_{I,{\rm dir}}^{\rm CIS(D)}
  +
  \Delta_{I,{\rm ind}}^{\rm CIS(D)},
  \label{eq:cisd-total-si}
\end{equation}
and the learned double-hybrid scaling is applied only through
$a_c^\theta$ in Eq.~\ref{eq:pt2-cisd-corrected-root-si}. This makes
the semilocal/HF response kernel and the perturbative doubles
correction separate mathematical objects while keeping both
differentiable with respect to the neural functional parameters.

% ======================================================================
\subsection{Neural XC Functional Architecture}
\label{sec:neuralxc}

The neural XC functional in GradTDDFT follows the
\emph{coefficient-basis} paradigm used by GradDFT-style neural
functionals. A pointwise neural network maps local descriptors to
mixing coefficients, and those coefficients multiply a fixed set of
energy-density channels. At grid point $g$,
\begin{equation}
  e_{\rm xc}^{\theta}(g)
  =
  \sum_{k=1}^{n_{\rm sl}}
  c_k^{\theta}(\mathbf{z}_g)\,e_k^{\rm sl}(g)
  + \mathbb{1}_{\rm pt2}\,
    c_{\rm pt2}^{\theta}(\mathbf{z}_g)\,p_g
  + c_{\rm hf}^{\theta}(\mathbf{z}_g)\,h_g,
  \label{eq:neural-eloc}
\end{equation}
and
\begin{equation}
  E_{\rm xc}^{\theta}
  =
  \sum_g w_g e_{\rm xc}^{\theta}(g).
  \label{eq:totalexc}
\end{equation}
The channel densities in Eq.~\ref{eq:neural-eloc} are defined on
the same quadrature grid. A semilocal channel has the general
spin-resolved meta-GGA form
\begin{equation}
  E_k^{\rm sl}
  =
  \int e_k^{\rm sl}(\mathbf{r})\,d\mathbf{r},
  \qquad
  e_k^{\rm sl}(\mathbf{r})
  =
  \rho(\mathbf{r})\,
  \varepsilon_k^{\rm sl}
  \!\left(
    \rho_\alpha,\rho_\beta,
    \sigma_{\alpha\alpha},\sigma_{\alpha\beta},
    \sigma_{\beta\beta},
    \tau_\alpha,\tau_\beta
  \right)_{\mathbf{r}},
  \label{eq:semilocal-density}
\end{equation}
where
$\rho=\rho_\alpha+\rho_\beta$,
$\sigma_{\sigma\sigma'}=\nabla\rho_\sigma\cdot
\nabla\rho_{\sigma'}$, and $\tau_\sigma$ is the spin kinetic-energy
density. GGA channels are obtained by omitting $\tau_\sigma$, and
LDA channels by omitting both gradient and kinetic-energy variables.
On the grid, $E_k^{\rm sl}\simeq\sum_g w_g e_k^{\rm sl}(g)$.

The semilocal channel set is not restricted to the B3LYP-like
choice below. Through jax-xc,\cite{JAXXC} translated Libxc LDA,
GGA, and meta-GGA energy densities can be inserted as channels,
\begin{equation}
  \mathcal{B}_{\rm sl}
  =
  \left\{
    e_{\rm xc}^{F_1}(g),\ldots,e_{\rm xc}^{F_M}(g)
  \right\},
  \label{eq:jaxxc-basis}
\end{equation}
where each $F_m$ denotes a conventional semilocal functional form.
A fixed functional mixture is recovered by constant coefficients
$a_m$,
\begin{equation}
  e_{\rm xc}^{F}(g)=\sum_{m=1}^{M}a_m e_{\rm xc}^{F_m}(g),
  \label{eq:fixed-functional-mixture}
\end{equation}
whereas GradTDDFT replaces $a_m$ by learned pointwise fields
$c_m^\theta(\mathbf{z}_g)$. This gives neural-network versions of
different semilocal functional families without changing the SCF or
TDDFT response equations.

The HF channel is a projected local representation of the
exact-exchange energy, following the same local-energy-density
viewpoint used in doubly local double-hybrid functionals.\cite{KovacsDLDH2026}
In spin-orbital notation,
\begin{equation}
  E_x^{\rm HF}
  =
  -\frac{1}{2}
  \sum_{\sigma}\sum_{ij}^{\rm occ}
  n_{i\sigma}n_{j\sigma}
  \iint
  \frac{
    \phi_{i\sigma}(\mathbf{r})
    \phi_{j\sigma}(\mathbf{r})
    \phi_{j\sigma}(\mathbf{r}')
    \phi_{i\sigma}(\mathbf{r}')
  }
  {|\mathbf{r}-\mathbf{r}'|}
  \,d\mathbf{r}\,d\mathbf{r}'.
  \label{eq:hf-energy}
\end{equation}
The local descriptor used in Eq.~\ref{eq:neural-eloc} is the
exchange-hole-gauge density
\begin{equation}
  h_g
  =
  -\frac{1}{2}
  \sum_{\sigma}\sum_{ij}^{\rm occ}
  n_{i\sigma}n_{j\sigma}
  \phi_{i\sigma}(\mathbf{r}_g)
  \phi_{j\sigma}(\mathbf{r}_g)
  v_{ij\sigma}^{x}(\mathbf{r}_g),
  \qquad
  v_{ij\sigma}^{x}(\mathbf{r})
  =
  \int
  \frac{
    \phi_{j\sigma}(\mathbf{r}')
    \phi_{i\sigma}(\mathbf{r}')
  }
  {|\mathbf{r}-\mathbf{r}'|}
  \,d\mathbf{r}',
  \label{eq:hf-local-density}
\end{equation}
so that $\sum_g w_g h_g$ is the quadrature representation of
$E_x^{\rm HF}$ up to grid and projection error.

The optional PT2 channel uses the same local-channel construction.
It starts from the canonical second-order correlation expression
\begin{equation}
  E_c^{(2)}
  =
  \frac{1}{4}
  \sum_{ij}^{\rm occ}\sum_{ab}^{\rm virt}
  \frac{
    |\langle ij||ab\rangle|^2
  }
  {\epsilon_i+\epsilon_j-\epsilon_a-\epsilon_b},
  \qquad
  \langle ij||ab\rangle=(ij|ab)-(ij|ba),
  \label{eq:pt2-energy}
\end{equation}
where $i,j$ are occupied spin orbitals and $a,b$ are virtual spin
orbitals. GradTDDFT uses a local projected density $p_g$ satisfying
\begin{equation}
  E_c^{(2)}
  \simeq
  \sum_g w_g p_g,
  \qquad
  p_g
  =
  \frac{1}{4}
  \sum_{ijab}
  t_{ij}^{ab}\,\eta_{ij}^{ab}(\mathbf{r}_g),
  \qquad
  t_{ij}^{ab}
  =
  \frac{\langle ij||ab\rangle}
       {\epsilon_i+\epsilon_j-\epsilon_a-\epsilon_b},
  \label{eq:pt2-local-density}
\end{equation}
with $\eta_{ij}^{ab}$ denoting the chosen real-space projection of
the antisymmetrized pair interaction, normalized so that
$\sum_g w_g\eta_{ij}^{ab}(\mathbf{r}_g)\simeq\langle ij||ab\rangle$.
Thus the neural coefficient multiplies a pointwise PT2 energy
density while the underlying orbital-energy denominators and pair
interactions retain their standard perturbative meaning.

The semilocal channels $e_k^{\rm sl}$ are evaluated by the
JAX/libxc-compatible backend. The default B3LYP-like basis is
\begin{equation}
  \{e_x^{\rm LDA}, e_x^{\rm B88},
    e_c^{\rm VWN\mbox{-}RPA}, e_c^{\rm LYP}\},
  \label{eq:default-basis}
\end{equation}
with coefficient priors matching the B3LYP local decomposition and
a separate HF prior. If the PT2 channel is enabled, it is inserted
between the semilocal channels and the HF channel with a zero
initial prior.

\textbf{Input descriptors.} The model supports two pointwise input
families. The canonical input uses a spin-resolved descriptor:
\begin{equation}
  \mathbf{z}_g^{\rm can}
  =
  [
  \rho_{\alpha},\rho_{\beta},
  \sigma,\sigma_{\alpha\alpha},\sigma_{\beta\beta},
  \tau_{\alpha},\tau_{\beta},
  h_{\alpha}^{\omega_1},\ldots,h_{\alpha}^{\omega_m},
  h_{\beta}^{\omega_1},\ldots,h_{\beta}^{\omega_m},
  p
  ]_g,
  \label{eq:canonical-input}
\end{equation}
where $\sigma_{\alpha\beta}=\nabla\rho_\alpha\cdot\nabla\rho_\beta$,
$\tau$ is the kinetic energy density, $h^\omega$ are local
range-separated HF descriptors, and $p$ is the optional PT2 local
descriptor. The default implementation uses $m=2$ HF descriptor
channels when local HF features are available. The enhanced input
augments these variables with $\rho$, $\log(1+\rho)$, $\sqrt{\rho}$,
the full set of spin-gradient invariants, total $\tau$, and a
semilocal energy-density descriptor. These descriptors are all
computed on the PySCF quadrature grid but enter the training graph as
JAX arrays.

\textbf{Pointwise mixing network.} Let
$\mathbf{a}_g^{\theta}=N_\theta(\mathbf{z}_g)$ be the raw network
output. The final channel coefficients are bounded as
\begin{equation}
  c_g^{\theta} = s\,\operatorname{sigmoid}
  \left(\frac{\mathbf{a}_g^{\theta}}{s}\right),
  \label{eq:coefficient-transform}
\end{equation}
with scale $s=2$ by default. Semilocal coefficients are then clipped
to the configured kernel range, while HF and PT2 heads are clipped to
the unit interval before they are used as mixing fields. The output
dimension is
\begin{equation}
  n_{\rm out}=n_{\rm sl}+1+\mathbb{1}_{\rm pt2},
  \label{eq:output-dim}
\end{equation}
corresponding to semilocal channels, optional PT2, and HF.

GradTDDFT provides two MLP architectures. The default residual
network applies
\begin{equation}
  \mathbf{x}_0=\tanh\!\left(W_0\log(|\mathbf{z}|+\epsilon)+b_0\right),
  \qquad
  \mathbf{x}_{\ell+1}
  =
  \varphi\!\left(
    {\rm LayerNorm}(W_\ell\mathbf{x}_{\ell}+\mathbf{b}_\ell
    +\Pi_\ell\mathbf{x}_{\ell})
  \right),
  \label{eq:residual-mlp}
\end{equation}
where $\Pi_\ell$ is either the identity or a learned projection when
layer widths change. The default hidden width is six layers of 256
units, and the default residual-block activation is ELU. The lighter
simple MLP instead applies the smooth transform
$\frac{1}{2}\log(\mathbf{z}^2+\epsilon^2)$ followed by dense layers
and a configurable activation. Both networks are implemented in
JAX/Flax.\cite{Flax}

\textbf{HF and PT2 heads.} The HF channel $h_g$ is a projected local
exact-exchange contribution and the PT2 channel $p_g$ is a projected
second-order perturbative local contribution. The HF field produces
both a local coefficient and an effective scalar exchange fraction
\begin{equation}
  \alpha_{\rm eff}^{\theta}
  =
  \frac{\sum_g w_g\,\rho_g\,c_{\rm hf}^{\theta}(\mathbf{z}_g)}
       {\sum_g w_g\,\rho_g},
  \label{eq:alpha-eff}
\end{equation}
which enters Eq.~\ref{eq:hfkernel}. In the excited-state response,
the local HF descriptor field is treated as a fixed coefficient input
for the local Hessian, and the nonlocal TDDFT response uses the
scalar exchange fraction in Eq.~\ref{eq:alpha-eff}. The PT2 head
defines the scalar coefficient in Eq.~\ref{eq:pt2-scalar-correlation-si}
for the CIS(D)-type excited-state correction.

% ======================================================================
\subsection{Differentiable Objective Terms}
\label{sec:training}

The differentiable formulation supports several classes of scalar
objectives. Rather than prescribing a fixed training sequence, we
write the objective as a weighted combination of ground-state,
and excited-state terms:
\begin{equation}
  \mathcal{J}(\theta)
  =
  \lambda_{\rm GS}\mathcal{L}_{\rm GS}(\theta)
  +
  \lambda_{\rm ES}\mathcal{L}_{\rm ES}(\theta),
  \label{eq:composite-loss}
\end{equation}
where the weights determine which physical constraints are active in
a given experiment.

\textbf{Ground-state terms.} Ground-state reference data can enter
through total energies, densities, or both:
\begin{equation}
  \mathcal{L}_{\rm GS} =
    w_E \cdot \text{MSE}(E^\theta, E^{\rm ref})
    + w_\rho \cdot \text{MAE}(\rho^\theta, \rho^{\rm ref}),
  \label{eq:gsloss}
\end{equation}
where $E^\theta$ is the SCF total energy obtained with the neural
functional and $\rho^\theta$ is the corresponding self-consistent
density. Depending on the chosen differentiable mode in
Eqs.~\ref{eq:fd-objective}--\ref{eq:implicit-scf-gradient}, the
density may be fixed, propagated through a finite SCF trajectory, or
treated as the solution of an implicit fixed-point problem.

\textbf{Excited-state terms.} Excited-state reference data enter
through vertical excitation energies and oscillator strengths:
\begin{equation}
  \mathcal{L}_{\rm ES} =
    w_\Omega \cdot \text{MAE}(\Omega_I^\theta, \Omega_I^{\rm ref})
    + w_f \cdot \text{MAE}(f_I^\theta, f_I^{\rm ref}),
  \label{eq:esloss}
\end{equation}
where $\Omega_I$ and $f_I$ are the excitation energies and
oscillator strengths for selected excited states $I$. The derivative
of this loss includes the dependence of the local energy density
$e_{\rm xc}^{\theta}$, XC potential components, local
$f_{\rm xc}$ tensor, orbital Hessian matrices, Casida eigenvectors,
and oscillator strengths on the neural parameters.
